\section{Wave propagation in three dimensions}\label{sec:waves3d}
\subsection{Plane waves in an isotropic generalized thermoelastic solid}
In an unbounded isotropic solid, with any heat-flux law of the template, \eqref{eq:helm} becomes
\begin{equation}
\rho\ddot\phi=(\lambda+2\mu)\nabla^2\phi-\beta\theta,\quad
s\rho c\,\hat\theta+T_0\beta s\nabla^2\hat\phi=G(s)\nabla^2\hat\theta,\quad
\rho\ddot{\bm\psi}=\mu\nabla^2\bm\psi,
\end{equation}
and a plane wave in direction $\bm n$, $\propto\exp[\ii(k\,\bm n\!\cdot\!\bx-\omega t)]$, obeys the same characteristic equation \eqref{eq:disp} for the dilatational--thermal system. The analysis of Section~\ref{sec:longitudinal} therefore holds for propagation in any direction of an isotropic body. Shear waves keep the speed $c_2$ and feel no heat-flux law, whether the model is a member of the template, fractional or memory-dependent. For general anisotropy the quasi-longitudinal, quasi-shear and thermal branches couple through $\beta_{ij}$ and $k_{ij}$, and the characteristic equation becomes a polynomial of higher degree; Dhaliwal and Sherief's formulation \citep{Dhaliwal1980} is the standard reference, and for numerical dispersion the Stroh formalism \citep{Stroh1958,Stroh1962} is convenient.

\paragraph{Counting branches.} In an isotropic medium a bulk plane-wave problem has $2$ branches from the dilatational--thermal pair $+\,2$ shear polarisations when $G$ is rational of low degree; the polynomial in $K$ has degree $2$ only for a $G$ of the template's kind. For the second-order lag version (ii) \eqref{eq:dpl2}, or higher-order models, the number of thermal-type branches increases. Published dispersion curves of such models show several further branches, and genuine physical branches then have to be separated from spurious roots created by Taylor truncation.

\subsection{Reflection at a free surface: structure of the problem}
A plane dilatational wave incident on a plane free surface in an isotropic thermoelastic half-space generates reflected waves of three kinds: a dilatational (elastic-type) wave, a thermal-type wave (evanescent or strongly damped), and a shear-vertical (SV) wave. Boundary conditions are the traction-free conditions together with a thermal condition, either \emph{insulated} (zero normal flux, $\partial_n\theta=0$ for any $G$) or \emph{isothermal} ($\theta=0$). A $3\times3$ linear system gives the amplitude ratios, and they have been computed for many models. The comparison of insulated and isothermal boundaries in generalized theories is a classical exercise \citep{Das2008}, and extensions with nonlocal, micropolar, piezoelectric and rotational effects are numerous \citep{Kalkal2020,Singh2020b,Yadav2023,Gupta2024a}. We do not tabulate those results. Section~\ref{sec:surface} gives instead the structure of the free-surface problem used for Rayleigh waves, which is the same determinant, and Section~\ref{sec:interfaces} treats the interface problem with a verified energy check.

\begin{figure}[t!]\centering\small
% fig4_reflection.tex -- Half-space reflection geometry (s08_waves3d)
\begin{tikzpicture}[x=1.0cm,y=1.0cm,>=Stealth]
  \tikzset{
    node/.style={draw,rounded corners=1pt,fill=blue!7,align=center,font=\footnotesize,inner sep=2pt,text width=2.2cm},
    arr/.style={->,thick,shorten >=1pt,shorten <=1pt},
    lbl/.style={font=\footnotesize,text=blue!60!black},
    medium/.style={fill=blue!10},
    surfaceline/.style={thick},
  }
  % Half-space region y>0 shaded
  \fill[medium] (-3.5,0) -- (3.5,0) -- (3.5,3.5) -- (-3.5,3.5) -- cycle;
  % Free surface line
  \draw[surfaceline] (-3.5,0) -- (3.5,0);
  \node[below=2pt] at (0,0) {free surface $y=0$};
  % Axes
  \draw[->,thick] (-3.5,0) -- (3.5,0) node[right] {$x$};
  \draw[->,thick] (0,0) -- (0,3.5) node[above] {$y$};
  % Incident P wave from upper left
  \draw[arr,very thick] (-2.5,2.0) -- (-0.7,0.6);
  \node[above left,lbl] at (-2.5,2.0) {incident P};
  % Reflected P wave
  \draw[arr,very thick] (-0.7,0.6) -- (1.5,2.2);
  \node[above right,lbl] at (1.5,2.2) {reflected P};
  % Reflected SV wave
  \draw[arr,very thick] (-0.7,0.6) -- (1.5,-0.8);
  \node[below right,lbl] at (1.5,-0.8) {reflected SV};
  % Thermal-type (evanescent) wave going down
  \draw[arr,dashed,very thick] (-0.7,0.6) -- (0.5,-2.0);
  \node[below,lbl] at (0.5,-2.0) {thermal (evanescent)};
  % Region label
  \node[above] at (0,3.2) {half-space $y>0$};
  % Boundary condition labels
  \node[below=6pt,align=center] at (-1.2,0) {insulated\\(\\partial_n\\theta=0)\\\\\nor isothermal\\($\\theta=0$)};
\end{tikzpicture}

\caption{Geometry after \citep{Das2008}: reflection at a free thermoelastic surface -- an incident P wave generates reflected P, reflected SV and a thermal-type (evanescent/damped) wave; thermal condition insulated or isothermal.}\label{fig:reflection-geo}
\end{figure}

\subsection{Waves in plates, rods and layered structures}
Lateral boundary conditions couple the thermal and elastic fields in bounded waveguides and produce families of guided modes: Lamb modes in plates \citep{Lamb1917} and Pochhammer--Chree modes in rods. The thermoelastic version adds one family of thermal modes per thermal condition, and the thermal coupling perturbs the elastic dispersion curves as in \Cref{fig:disp}, with small speed corrections and damping that peaks where the thermal relaxation frequency of the cross-section matches the wave frequency. For thin plates and beams the flexural damping reduces to Zener's classical result (Section~\ref{sec:ted}). The elastic waveguide problem has textbooks \citep{Achenbach1973,Auld1973}; for the thermoelastic waveguide with generalized laws the literature consists mainly of specific configurations \citep{Biswas2020,Abouelregal2023}.

\begin{summarybox}[Summary]
In an unbounded isotropic body plane waves split into a shear branch (speed $c_2$, unaffected by heat) and a coupled dilatational--thermal system governed by \eqref{eq:disp}. Reflection at a free surface generates three reflected waves (P, SV, thermal-type), with amplitude ratios from a $3\times3$ system. In waveguides each thermal condition adds a family of thermal modes.
\end{summarybox}
